\documentclass[
  fontsize=12pt  %
 ,english        % 
 ,headinclude    %
 ,headsepline    % line between head an document text
%,BCOR=12mm      % 
]{scrbook}       % twosided, A4 paper

\usepackage[T1]{fontenc}

\usepackage{polski}
\usepackage[utf8]{inputenc}
\usepackage[french,polutonikogreek,polish]{babel}

%\geometry{verbose,a4paper,tmargin=3cm,bmargin=2cm,lmargin=2cm,rmargin=2cm}


\usepackage{blindtext} % provides blindtext with sectioning

\usepackage{scrpage2}  % header and footer for KOMA-Script

\usepackage{graphicx}

\clearscrheadfoot                 % deletes header/footer
\pagestyle{scrheadings}           % use following definitions for header/footer
% definitions/configuration for the header
\rehead[]{\Large \textbf{BitAlgo Start}}        % equal page, right position (inner) 
\lohead[]{\Large \textbf{BitAlgo Start}}        % odd   page, left  position (inner) 
\cehead[]{Zadanie V:\\Kolejka}
\cohead[]{Zadanie V:\\Kolejka}
\lehead[]{\includegraphics[width=15mm]{logo.png}} % equal page, left (outer) position
\rohead[]{\includegraphics[width=15mm]{logo.png}}
% definitions/configuration for the footer
\cofoot[\pagemark]{\pagemark}     % odd   page, center position

\begin{document}
\vspace{50 mm}
\hspace{50 mm}
\newline

\par{\Large \textbf{Zadanie V: Kolejka}} \\ \\
Twoim zadaniem jest efektywna implementacja interfejsu kolejki (FIFO) z ograniczoną pojemnością. Interfejs zawiera operację sprawdzania wartości i usuwania pierwszego elementu, dodawania elementu na końcu. Zakładamy, że na samym początku kolejka jest pusta.
\\
Uwaga! Nie korzystaj z kontenera std::queue lub std::deque. Twoja implementacja powinna używać tablicy a nie listy. Wykorzystaj ideę cyklicznego bufora.
\\ \\
\par{\Large \textbf{Format wejścia}} \\ \\
Pierwsza linia wejścia dwie liczby całkowite $n$ i $k$, $0 \leq n, k \leq 10^6$ - ilość poleceń. $n$ definiuje pojemność kolejki, tzn. maksymalną ilość elementów. W kolejnych $k$ liniach znajduje się jedno polecenie z poniższych:
\begin{itemize}
  \item $PUSH x$ - dodaje element $x$ na koniec kolejki,
  \item $POP$ - usuwa pierwszy element z kolejki,
  \item $FRONT$ - zwraca wartość pierwszego elementu w kolejce.
\end{itemize}

Wartości elementów są ograniczone: $0 \leq x \leq 10^6$.
\\ \\
\par{\Large \textbf{Format wyjścia}} \\ \\
Na wyjściu dla każdej komendy powinien zostać wypisany wynik jej działania:
\begin{itemize}
  \item $PUSH$ - napis $OK$ lub $FULL$ jeśli kolejka jest pełna,
  \item $POP$ - napis $OK$ lub $EMPTY$ jeśli kolejka jest pusta,
  \item $FRONT$ - wartość pierwszego elementu w kolejce lub napis $NONE$ jeśli kolejka jest pusta.
\end{itemize}

~\\
\vspace{50 mm}
\hspace{50 mm}
\newline
\newline

\par{\Large \textbf{Przykład}} \\ \\
\begin{tabular}{ p{7cm} p{7cm} }

  Dla danych wejściowych: \hspace{40mm}& Poprawną odpowiedzią jest \\
& \\

% define input here
3 11 \newline
FRONT \newline
PUSH 4 \newline
PUSH 3 \newline
PUSH 9 \newline
PUSH 8 \newline
FRONT \newline
POP \newline
POP \newline
FRONT \newline
POP \newline
POP \newline
&   
% define output here
NONE \newline
OK \newline
OK \newline
OK \newline
FULL \newline
4 \newline
OK \newline
OK \newline
9 \newline
OK \newline
EMPTY \newline

\\

\end{tabular}
\end{document}